\chapter{Model calibration}
\label{cap:swe_calibration}

\section{HS data}

The snow depth dataset used to calibrate the degree-day model described in the previous chapter was created through a data rescue activity. It combines 
HS time series from \cite{Matiu_HS} with data provided by national and regional meteorological and hydrological services, including ARPA Piemonte, Centro 
Funzionale della Valle d'Aosta, Meteoswiss, WSL Institute for Snow and Avalanche Research SLF, ARPA Lombardia, Meteotrentino, ARPA Veneto, and the Austrian 
Hydrological Service \parencite{hs_data}. The final dataset includes 1052 stations, mostly automatic, distributed across the entire Alpine Arc, as shown 
in Figure \ref{fig:hs_series_origin}.
\begin{figure}[H]
    \centering
    \includegraphics[width=\linewidth]{swe_calibration/hs_series_origin.jpg}
    \caption{Spatial distribution of the stations that made up the original dataset. Each station is marked with the symbol associated with the 
    data source from which the data were obtained. Image from \parencite{hs_data}.}
    \label{fig:hs_series_origin}
\end{figure}

The original HS dataset, however, was affected by several issues that needed to be addressed before it could be considered suitable for 
further analysis. The first reliability check that was done focused on assessing the accuracy of the station metadata, namely longitude, latitude 
and elevation. In particular, the ALOS AW3D30 digital elevation model was used to compare the reported elevation of each station with that of the DEM 
pixel closest to its coordinates. A compatibility threshold of $\SI{12}{\meter}$ was adopted: although this criterion is rather stringent given the 
complex topography of the Alpine Arc, it was considered appropriate given the high spatial resolution of the DEM, namely $\SI{30}{\meter}$. With this 
threshold, more than half of the stations were classified as unreliable, as show in Figure \ref{fig:station_distribution_original}.
\begin{figure}[H]
    \centering
    \includegraphics[width=\linewidth]{swe_calibration/station_distribution_original.png}
    \caption{(a) Spatial distribution of reliable and unreliable stations. (b) Frequency histogram of the elevation differences between the reported 
    station elevations and those derived from the ALOS AW3D30 DEM. In both panels, green indicates reliable stations, i.e. stations for which the 
    absolute difference between the reported and DEM-derived elevations does not exceed $\SI{12}{\meter}$, while red is used for stations that did 
    not satisfy this criterion.}
    \label{fig:station_distribution_original}
\end{figure}
that needed to be taken care of before labelling it as usable. The two main issues, presented in 
Figure , were missing data (i.e. reading gaps) across the time series and summer values that weer not going to zero when clearly the ground should be bare. 
In particular, the latter numerical error came in two form, namely as a homogenoeus bias across the whole year, or as increasing values as the season went on. 
\begin{figure}[H]
    \centering
    \includegraphics[width=\linewidth]{swe_calibration/comparison_HS_MODIS_original.png}
    \caption{(a) Spatial distribution of station }
    \label{fig:comparison_HS_MODIS_original}
\end{figure}

\subsection{Position correction}

\begin{figure}[H]
    \centering
    \includegraphics[width=\linewidth]{swe_calibration/station_compatibility_before_after_QC.png}
    \caption{(A) Compatibility between declared station elevation and closest DEM point before checking station coordinates. (B) Compatibility 
    between declared station elevation and closest DEM point after checking station coordinates. In both cases the ALOS AW3D30 global DEM was used: 
    stations are labelled as compatible (OK) if the elevation difference is smaller than $\SI{12}{\meter}$, namely 2/5 of pixel size.}
    \label{fig:station_compatibility_before_after_QC}
\end{figure}

\begin{figure}[H]
    \centering
    \includegraphics[width=\linewidth]{swe_calibration/elevation_distribution_stations.png}
    \caption{}
    \label{fig:elevation_distribution_stations}
\end{figure}

\subsection{HS series correction}

\begin{figure}[H]
    \centering
    \includegraphics[width=\linewidth]{swe_calibration/comparison_HS_MODIS_corrected.png}
    \caption{}
    \label{fig:comparison_HS_MODIS_corrected}
\end{figure}

